% Part: incompleteness
% Chapter: incompleteness-provability
% Section: provability-conditions

\documentclass[../../../include/open-logic-section]{subfiles}

\begin{document}

\olfileid{inc}{inp}{prc}

\olsection{Sarat \usetoken{S}{derivability} kanggo $\Th{PA}$}

Aritmetika Peano, utawa $\Th{PA}$, iku teori kang ngembangake
$\Th{Q}$ nganggo aksioma induksi kanggo saben !!{formula}s.
Kanthi tembung liya, ing $\Th{Q}$ ditambahake aksioma awujud
\[
(!A(0) \land \lforall[x][(!A(x) \lif !A(x'))]) \lif \lforall[x][!A(x)]
\]
kanggo saben !!{formula}~$!A$. Gatekna yen iki sejatine sawijining
\emph{skema}, yaiku aksioma kang cacahe tanpa wates (lan jebul
$\Th{PA}$ {\em ora} bisa diaksiomatisasi kanthi winates). Nanging
amarga kita bisa kanthi efektif mutusake apa sawijining rerangkening
simbol iku conto saka aksioma induksi utawa ora, himpunan aksioma
kanggo $\Th{PA}$ komputabel. $\Th{PA}$ iku teori kang luwih kuwat
tinimbang~$\Th{Q}$. Upamane, kanthi induksi lumrah kita gampang
mbuktekake yen panjumlahan lan pingan komutatif. Malah, akeh
pambuktian finitari ing teori bilangan lan kombinatorika bisa
ditindakake ing~$\Th{PA}$.

Amarga $\Th{PA}$ bisa diaksiomatisasi kanthi komputabel, predikat
!!{derivability} $\Prf[\Th{PA}](x,y)$ komputabel lan mula bisa
direpresentasikake ing~$\Th{Q}$ (lan uga ing~$\Th{PA}$). Kaya
sadurunge, $\OPrf[\Th{PA}](x,y)$ kita gunakake kanggo nyebut formula
kang merepresentasikake relasi iki. Tetepna $\OProv[\Th{PA}](y)$
minangka formula $\lexists[x][\Prf[\Th{PA}](x,y)]$, kang sacara
intuitif nyatakake ``$y$ !!{derivable} saka aksioma-aksioma
$\Th{PA}$.'' Alesan kita mbutuhake luwih saka aksioma $\Th{Q}$
yaiku kita kudu ngerti yen teori kang digunakake cukup kuwat
kanggo !!{derive} sawetara fakta dhasar ngenani predikat
!!{derivability} iki. Nyatane, fakta kang dibutuhake yaiku:
\begin{enumerate}
\item[P1.] Yen $\Th{PA} \Proves !A$, mula $\Th{PA} \Proves
  \OProv[\Th{PA}](\gn{!A})$.
\item[P2.] Kanggo saben !!{formula}s $!A$ lan $!B$,
  \[
  \Th{PA} \Proves \OProv[\Th{PA}](\gn{!A \lif !B}) \lif
  (\OProv[\Th{PA}](\gn{!A}) \lif \OProv[\Th{PA}](\gn{!B})).
  \]
\item[P3.] Kanggo saben !!{formula}~$!A$,
  \[
  \Th{PA} \Proves \OProv[\Th{PA}](\gn{!A})
  \lif \OProv[\Th{PA}](\gn{\OProv[\Th{PA}](\gn{!A})}).
  \]
\end{enumerate}
Kanggo mbuktekake yen telung sipat iki bener, kita kudu njlentrehake
!!{formula} $\OProv[\Th{PA}](y)$ kanthi tliti lan nggunakake aksioma
$\Th{PA}$ kanggo njlentrehake !!{derivation}s formal kang relevan.
Sarat (1) lan~(2) gampang; sarat~(3) kang sejatine mbutuhake
pakaryan. (Coba pikirna pakaryan apa kang dibutuhake \dots)
Rincian kasebut bakal dawa lan ora narik kawigaten, mula ing kene
kita njaluk sampeyan nampa tanpa bukti yen $\Th{PA}$ nduweni telung
sipat ing ndhuwur. Pilihan $\OProv[\Th{PA}](y)$ kang lumrah uga
bakal nyukupi
\begin{enumerate}
\item[P4.] Yen $\Th{PA} \Proves \OProv[\Th{PA}](\gn{!A})$, mula
  $\Th{PA} \Proves !A$.
\end{enumerate}
Nanging fakta iki ora bakal dibutuhake.

\begin{digress}
Minangka cathetan, G\"odel uga rada kesed kaya kita saiki. Ing
pungkasan makalahe taun 1931, dheweke nggambarake garis gedhe bukti
teorema ora jangkep kapindho lan janji bakal njlentrehake rinciane
ing makalah sabanjure. Nanging dheweke ora nate nindakake iku;
amarga saben wong kang mangerteni argumene percaya yen bukti iku
bisa dirampungake, dheweke ora perlu ngisi rinciane.
\end{digress}

\end{document}
